\section{Introduction} % 1.5 page

% LLMs
% The rapid growth of Large Language Model (LLM) serving drives new computational requirements. 
Since the development of the Transformer architecture~\cite{vaswani2017attention, kaplan2020scaling}, optimizing Large Language Model (LLM) inference for production has become a critical challenge~\cite{2025taming}. Modern serving frameworks have two main objectives: maintaining high aggregate throughput for cost efficiency, and meeting latency Service-Level Objectives~\cite{agrawal2024sarathi}. In highly interactive scenarios, the limiting factor becomes single-batch user input latency. Since GPUs have traditionally been designed as throughput-oriented devices, minimizing latency for a single-batch request introduces new system-level requirements.

% Multi-GPU systems
Correspondingly, to support the rapid scaling of LLM workloads, modern multi-GPU hardware systems have evolved into tightly coupled architectures, calling for novel software paradigm to fully exploit them. Systems such as the NVIDIA GB200 NVL72~\cite{nvidia2024gb200nvl72} utilize NVLink and NVSwitch~\cite{nvidia2024nvlink} interconnects not only to provide high peer-to-peer bandwidth, but also to enable the multi-GPU cluster to function with a unified shared-memory space. Although recent advancements have introduced sophisticated serving frameworks to maximize multi-GPU deployment efficiency~\cite{kwon2023vllm, zheng2023sglang, llmd2025, aibrix2025, nvidia2025dynamo}, as well as efforts to build systematic abstractions for programming multi-GPU workloads~\cite{sul2025parallelkittens, heldens2022lightning}, there remains untapped potential for novel execution paradigms that natively exploit these tightly coupled clusters as holistic shared-memory systems.

At the current stage, the standard paradigm for large-scale LLM deployment on multi-GPU systems relies on hybrid parallelism strategies~\cite{smith2022megatron, zheng2022alpa, narayanan2021efficient, qi2025synergistic}, primarily combining Pipeline Parallelism (PP)~\cite{huang2019gpipe, narayanan2019pipedream} and Tensor Parallelism (TP)~\cite{shoeybi2019megatron}. While emerging techniques such as Expert Parallelism (EP)~\cite{gale2023megablocks} and disaggregated serving~\cite{zhong2024distserve} offer further optimizations, they are highly workload-specific. Therefore, PP and TP remain the universal baselines. PP mainly focuses on improving overall serving throughput by operating at the inter-layer level to distribute transformer blocks across devices. TP provides both throughput improvement and latency reduction by spatially sharding computations at the operator level. However, TP introduces additional collective communication (e.g., AllReduce) to resolve data dependencies~\cite{zheng2022alpa}, establishing a hard ceiling on latency optimization.

Between inter-layer PP and operator-level TP lies an opportunity for further latency reduction: the \textit{intra-layer, inter-operator} execution space. Existing efforts accelerate this space primarily through temporal optimizations to improve efficiency on single device, such as kernel fusion~\cite{dao2022flashattention} or mega-kernels~\cite{aminabadi2022deepspeed, wu2024mirage,cheng2025mpk}. However, these require complex compiler toolchains or rigid rewrites~\cite{cheng2025mpk}. Alternatively, localized micro-batching~\cite{agrawal2023sarathi, zhong2024distserve} offers finer-grained pipelining, but introduces pipeline bubbles and degrades kernel efficiency at small chunk sizes. These limitations call for more efficient inter-operator spatial scaling techniques.

To address the dual requirements of single-batch latency optimization for GPU workloads and the need for new software paradigms for shared-memory multi-GPU systems, we propose CTA-pipelining, a novel latency-oriented spatial scaling method. Operating at the inter-operator level, it leverages the unified NVLink memory domain to enable simultaneous execution of data-dependent kernels across GPUs, via dynamic spatial pipelining at the Cooperative Thread Array (CTA) granularity. The designed protocol is minimally invasive to existing GPU kernel implementation, by relying only on the addition of prologue and epilogue code snippets. This preserves the potential for automated integration across a wide variety of workloads.

% Unlike traditional micro-batch chunk pipelining, our technique avoids fragmenting kernels into small, static chunks, thereby preserving kernel efficiency while maximizing cross-GPU parallelism. This design also ensures broad generalizability across a diverse range of GPU workloads, making it suitable as a general spatial scaling method.

As an initial demonstration, we implement prototype and perform analysis using multi-layer general matrix-matrix multiplication (GEMM), a critical GPU workload, which also represents the multilayer perceptron (MLP) layers in transformer-based LLMs. Our implementation and evaluation are built upon the state-of-the-art NVIDIA libraries including CUTLASS~\cite{cutlass}, cuBLAS~\cite{cublas}, and NCCL~\cite{nccl}, on cutting-edge hardware including 8-GPU H200 NVLink system and 8-GPU B200 NVLink system.

We evaluate CTA-pipelining via two angles: validating the fundamental mechanism and demonstrating its broader scaling capabilities. First, to validate the protocol's overhead, we analyze the integration with both classical and warp-specialized multi-stage persistent CUTLASS kernels. Results show the overhead is minimal, and can even be mostly hidden within warp-specialized executions. Second, to establish CTA-pipelining as a generalized spatial scaling paradigm, we evaluate CTA-pipelining against traditional micro-batch chunk pipelining (achieving up to 31.8\% latency reduction) and Tensor Parallelism (achieving up to 29.6\% latency reduction), on setups representing MLP operations. We further demonstrate that it can be combined with TP as an orthogonal spatial scaling dimension, further pushing the latency scaling limit by providing benefit on both computation and communication.

% Contribution
The contributions of this paper are summarized as follows:
\begin{itemize}
    \item We propose CTA-pipelining, a novel spatial scaling paradigm for multi-GPU shared-memory systems, aimed at optimizing single-batch latency for multi-GPU workloads such as LLM inference.
    \item We demonstrate a prototype implementation of CTA-pipelining using the state-of-the-art GEMM library, NVIDIA CUTLASS, on various styles of GPU kernels.
    \item We perform analysis on basic protocol overhead, and identify the benefit of integrating with warp-specialized multi-stage persistent kernels.
    \item We compare CTA-pipelining against traditional micro-batch chunk pipelining, yielding latency reductions of up to 31.8\% on multi-layer GEMMs. 
    \item We compare CTA-pipelining against Tensor Parallelism on multilayer perceptron setup, providing latency reductions of up to 29.6\%. 
    \item We show how CTA-pipelining serves as an orthogonal spatial scaling method to Tensor Parallelism, offering benefit on both computation and communication, further pushing the latency optimization frontier. 
    \item We hint potential hardware evolutions that could further benefit this execution model.
\end{itemize}

% The remainder of this paper is organized as follows. Section 2 introduces the background and existing work necessary to understand our approach, including related GPU execution models and current scaling methods for transformer-based LLMs. Section 3 details the design and implementation of CTA-pipelining, covering its general architecture and suitable use cases. Section 4 presents our experimental evaluation of workload performance, by analyzing protocol overhead across different kernel implementation styles, comparing our approach against traditional micro-batching, and evaluating the combination of CTA-pipelining with TP against a pure TP baseline. Section 5 provides discussions on technical decisions made during development, potential use cases, possible improvements, and how this work can guide future hardware evolution. Finally, the paper is concluded in Section 6.

